\chapter{Calcolo derivate}
\section{Regole di Derivazione Fondamentali}
Per calcolare le derivate senza ricorrere alla definizione di limite, utilizziamo le regole di derivazione:

\begin{itemize}
    \item \textbf{Derivata di una costante:} $D[c] = 0$
    \item \textbf{Derivata di una potenza:} $D[x^n] = n x^{n-1}$
    \item \textbf{Derivata di una somma:} $D[f(x) + g(x)] = f'(x) + g'(x)$
    \item \textbf{Regola del prodotto:} $D[f(x)g(x)] = f'(x)g(x) + f(x)g'(x)$
    \item \textbf{Regola del quoziente:} $D\left[\frac{f(x)}{g(x)}\right] = \frac{f'(x)g(x) - f(x)g'(x)}{[g(x)]^2}$
    \item \textbf{Regola della catena (funzione composta):} $D[f(g(x))] = f'(g(x)) \cdot g'(x)$
\end{itemize}

\section{Derivate di Funzioni Elementari}
\begin{table}[h!]
\centering
\begin{tabular}{|l|l|}
\hline
\textbf{Funzione $f(x)$} & \textbf{Derivata $f'(x)$} \\ \hline
$\sin(x)$ & $\cos(x)$ \\ \hline
$\cos(x)$ & $-\sin(x)$ \\ \hline
$e^x$ & $e^x$ \\ \hline
$\ln(x)$ & $1/x$ \\ \hline
\end{tabular}
\end{table}

